\documentclass[preprint,12pt]{elsarticle}

\usepackage{amsmath}
\usepackage{amssymb}
\usepackage{hyperref}

\begin{document}

\begin{frontmatter}

\title{A Python Library for the Theoretical-Numerical Search and Localization of Hidden Attractors}

\author[inaoe]{Maria Fernanda Moreno Lopez}
\affiliation[inaoe]{organization={Instituto Nacional de Astrofisica, Optica y Electronica (INAOE)}}

\begin{abstract}
This manuscript skeleton documents the software scope and reproducibility structure for a future Computer Physics Communications submission. The package, named \texttt{hidden-attractors-fo}, organizes workflows for theoretical-numerical search, localization, reproduction, audit, and conservative classification of candidate hidden attractors in integer- and commensurate fractional-order Chua/Lur'e systems. The current preparation does not introduce new numerical results. Scientific claims remain bounded by the promoted validation artifacts in \texttt{version\_2/validation/}, the finite-time numerical contract, and the tested-neighborhood hiddenness language used by the project.
\end{abstract}

\begin{keyword}
hidden attractors \sep fractional-order systems \sep integer-order systems \sep Chua circuit \sep Caputo derivative \sep numerical validation \sep scientific software
\end{keyword}

\end{frontmatter}

\section*{Program Summary}

\noindent\textbf{Program title:} hidden-attractors-fo

\noindent\textbf{Archived title:} A Python Library for the Theoretical-Numerical Search and Localization of Hidden Attractors

\noindent\textbf{Repository:} \url{https://github.com/Xerkkun/Hidden-Attractors-Localization}

\noindent\textbf{Licensing provisions:} MIT

\noindent\textbf{Programming language:} Python, with optional C backends

\noindent\textbf{Supplementary material:} Archived software/material DOI: \url{https://doi.org/10.17605/OSF.IO/ZGK74}

\noindent\textbf{Nature of problem:} Theoretical-numerical search, localization, reproduction, audit, and conservative classification of candidate hidden attractors in integer- and commensurate fractional-order Chua/Lur'e systems with Caputo derivatives for the fractional case.

\noindent\textbf{Solution method:} Describing-function seeding, BDF-style algebraic continuation, integer-order integration paths, Caputo ABM/EFORK integration for fractional order, finite-time chaos diagnostics, and equilibrium-neighborhood tests around all equilibria.

\noindent\textbf{Additional comments including restrictions and unusual features:} The implementation targets scalar Lur'e systems and commensurate fractional order in the fractional workflow. Numerical neighborhood tests do not certify global mathematical hiddenness. The arctan Chua route is implemented algebraically but is not promoted as a validated hidden attractor.

\noindent\textbf{References:} See \texttt{references.bib}.


\section{Introduction}

Hidden attractors and the distinction between self-excited and hidden oscillations are central to the motivation for this package \cite{leonov_kuznetsov_hidden_2013,kuznetsov_hidden_attractors_2017}. The Chua circuit remains a reference setting for nonlinear circuit dynamics and chaos studies \cite{matsumoto_chua_1984}. This repository prepares a reproducible software layer for candidate localization and audit, not a mathematical proof engine.

\section{Scientific and numerical scope}

The active distribution is \texttt{version\_2/}. It supports integer-order and commensurate fractional-order workflows for scalar Lur'e-compatible Chua-type systems. Fractional workflows are written under Caputo-style memory assumptions and use explicit solver, horizon, step-size, and memory-policy contracts \cite{podlubny_fractional_1999}. The terminology used in reports is intentionally conservative: candidate hidden attractor, finite-time numerical evidence, tested neighborhoods, reproduced reference case, and not a global proof.

The Chua arctan route is implemented algebraically and is useful for traceability against the documented source case \cite{wu_chua_arctan_2023}. It remains pending full validation and is not promoted as a validated hidden attractor.

\section{Software architecture}

The package exposes a single public entry point, \texttt{hidden-attractors}. Core code lives in \texttt{version\_2/hidden\_attractors/}; configurations live under \texttt{version\_2/configs/}; promoted validation evidence lives under \texttt{version\_2/validation/}; promoted scientific figures belong in \texttt{version\_2/library\_figures/}; and local/regenerable outputs remain ignored under \texttt{version\_2/outputs/}, \texttt{version\_2/validation\_outputs/}, \texttt{version\_2/runs*/}, or \texttt{version\_2/figures/}. Editorial preparation material is separated under \texttt{paper/} and \texttt{version\_2/cpc\_submission/}.

\section{Numerical workflow}

The workflow combines describing-function or harmonic-balance style seed generation, BDF-style algebraic continuation, integer-order integration paths, Caputo ABM/EFORK integration for fractional order, and finite-time diagnostics \cite{diethelm_abm_fractional_2002,ghoreishi_ghaffari_saad_efork_2023,khalil_nonlinear_2002}. Chaos diagnostics such as the 0-1 test and Lyapunov-style estimates are treated as supporting evidence only \cite{gottwald_melbourne_2004,wolf_lyapunov_1985}. They do not replace equilibrium-neighborhood tests.

\section{Validation and reproducibility}

Promoted evidence is checked through the repository validation contract. The integer Chua case is treated as a reproduced reference/software validation case. Fractional non-smooth Chua evidence is reported using the project claims matrix and finite-time neighborhood tests; conclusions remain bounded to the recorded contract. Regenerable outputs are intentionally excluded from Git to keep the promoted evidence boundary auditable.

The current CPC preparation marks the freeze audit as \texttt{pending\_after\_cpc\_cleanup}. The last recorded freeze audit predates this CPC cleanup and must be regenerated before a final submission claim is made.

\section{Limitations}

The software does not certify global mathematical hiddenness. It records finite-time numerical evidence under tested neighborhoods, chosen integration horizons, solver tolerances, sampling radii, and memory policies. Describing-function, Nyquist, continuation, boundedness, and diagnostic plots do not prove hiddenness by themselves. Arctan Chua remains implemented algebraically but pending full validation.

\section{Availability and archival metadata}

The repository is available at \url{https://github.com/Xerkkun/Hidden-Attractors-Localization}. The current archived software/material DOI is \url{https://doi.org/10.17605/OSF.IO/ZGK74}. CPC Library metadata should be assigned only after acceptance.

\section{Conclusions}

This preparation organizes the package, evidence boundaries, metadata, and reproducibility checks for a future CPC submission. The manuscript remains a preparation template; it should be completed with regenerated freeze-audit results and final bibliographic metadata before submission.

\bibliographystyle{elsarticle-num}
\bibliography{references}

\end{document}
